\chapter{Self-Harm}

\section*{Definition}
Non-suicidal self-injury (NSSI) is the deliberate, self-inflicted destruction of body tissue without suicidal intent, such as cutting, burning, scratching, or hitting. It is distinguished from suicidal behavior by the absence of an intent to die.

\section*{What Happens in Your Body}
NSSI is thought to dysregulate the endogenous opioid and dopaminergic systems, with self-injury providing temporary relief from intense negative affect through the release of endogenous opioids and endorphins. Altered pain perception (higher pain tolerance during episodes) and deficits in emotion regulation circuitry (prefrontal--amygdala) are characteristic.

\section*{Causes}
\begin{itemize}
  \item \textbf{Psychiatric:} Borderline personality disorder, major depressive disorder, PTSD, eating disorders, substance use disorders, anxiety disorders, autism spectrum disorder
  \item \textbf{Psychological factors:} Emotional dysregulation, intense anger or shame, dissociation, self-criticism, alexithymia, poor distress tolerance
  \item \textbf{Developmental:} History of childhood abuse (physical, emotional, sexual), neglect, insecure attachment, bullying
  \item \textbf{Social:} Peer modeling (contagion effect, especially in adolescents), social media normalization, identity expression
  \item \textbf{Genetic:} Heritability ~40--60\%; associated with variants in serotonin and opioid receptor genes
\end{itemize}

\section*{When to See a Doctor}
See your doctor if:
\begin{itemize}
  \item Self-injury occurs with suicidal intent, or there is passive or active suicidal ideation
  \item Wounds are deep, infected, or require medical attention
  \item Frequency is escalating, or the individual feels unable to control the behavior
\end{itemize}

\section*{Related Symptoms}
\begin{itemize}
  \item Depression, suicidal ideation, emotional dysregulation, impulsivity, dissociation, eating disorders, substance use, low self-esteem, borderline personality disorder
\end{itemize}
\textbf{Important:} NSSI typically begins between ages 13 and 15 and is a strong risk factor for future suicide attempts, regardless of stated intent.

\section*{Self-Care / Home Management}
\begin{itemize}
  \item Use less harmful alternatives: snapping a rubber band on the wrist, holding ice cubes, drawing on skin with red marker
  \item Practice the STOP skill: Stop, Take a step back, Observe, Proceed mindfully
  \item Create a safety plan with a list of supportive contacts and distraction activities
  \item Reach out to a crisis text line or warm line when urge is strong
\end{itemize}

\section*{How Your Doctor Will Evaluate You}
\begin{itemize}
  \item Clinical interview distinguishing NSSI from suicidal behavior, assessing frequency, methods, triggers, and intent
  \item Assess for psychiatric comorbidities (borderline personality disorder, depression, PTSD, eating disorders)
  \item Structured scales: Self-Injurious Thoughts and Behaviors Interview (SITBI), Inventory of Statements About Self-Injury (ISAS)
  \item Suicidal ideation and risk assessment are mandatory; use Columbia-Suicide Severity Rating Scale (C-SSRS)
\end{itemize}

\section*{Treatment Options}
\begin{itemize}
  \item First-line psychotherapy: Dialectical behavior therapy (DBT) -- specifically designed for NSSI and BPD
  \item Other evidence-based therapies: cognitive-behavioral therapy, mentalization-based therapy (MBT), cognitive restructuring
  \item No FDA-approved pharmacotherapy for NSSI; treat underlying depression, anxiety, or PTSD with SSRIs
  \item Safety planning, means restriction, and family involvement for adolescents
  \item Hospitalization if imminently suicidal or NSSI is severe and uncontrolled
\end{itemize}

\section*{References}
\begin{thebibliography}{99}
\bibitem{self_harm:ref1} Hawton K, Saunders KEA, O'Connor RC. "Self-harm and suicide in adolescents." \textit{The Lancet}, 2012;379(9834):2373--2382.
\bibitem{self_harm:ref2} Nock MK. "Self-injury." \textit{Annual Review of Clinical Psychology}, 2010;6:339--363.
\bibitem{self_harm:ref3} Klonsky ED. "The functions of deliberate self-injury: A review of the evidence." \textit{Clinical Psychology Review}, 2007;27(2):226--239.
\bibitem{self_harm:ref4} Ougrin D, Tranah T, Stahl D, et al. "Therapeutic interventions for suicide attempts and self-harm in adolescents: A systematic review." \textit{Journal of the American Academy of Child \& Adolescent Psychiatry}, 2015;54(2):97--107.
\bibitem{self_harm:ref5} Kothgassner OD, Robinson K, Goreis A, et al. "Does treatment method matter? A meta-analysis of the effectiveness of interventions for non-suicidal self-injury." \textit{Journal of Consulting and Clinical Psychology}, 2020;88(12):1119--1131.
\bibitem{self_harm:ref6} American Academy of Child and Adolescent Psychiatry. "Practice parameter for the assessment and treatment of children and adolescents with suicidal behavior." \textit{Journal of the American Academy of Child \& Adolescent Psychiatry}, 2001;40(7 Suppl):24S--51S.
\end{thebibliography}
